\section{Thảo luận}
\label{sec:discussion}

\paragraph{Khung diễn giải theo hai câu hỏi nghiên cứu.}
Vì số liệu cuối cùng được điền tự động từ artifact, phần này trình bày khung diễn giải
đã định trước theo từng kịch bản kết quả, thay vì diễn giải hậu nghiệm. Với RQ1, nếu
AlignmentGraph-ISD giữ thứ hạng dẫn đầu trên cả bốn kích thước và
difference-in-differences cho thấy lợi thế lớn hơn ở kích thước nhỏ với khoảng tin cậy
loại trừ 0, kết quả ủng hộ luận điểm harness-over-model~\cite{cho2026harness}: cách tổ
chức môi trường thực thi bù một phần cho giới hạn của mô hình nền, và phần bù này lớn
dần khi mô hình yếu đi. Nếu lợi thế ổn định nhưng không nới rộng, kết luận hẹp hơn:
harness cải thiện chất lượng đồng đều nhưng không đặc biệt hữu ích hơn cho SLM. Nếu
lợi thế thu hẹp ở kích thước nhỏ, giả thuyết thiết kế bị thách thức---các cơ chế của
harness có thể đòi hỏi một năng lực mô hình tối thiểu để vận hành.
\TBD{chốt diễn giải theo kịch bản xảy ra khi có số thực}.

Với RQ2, các chỉ số alignment cho phép tách hai khả năng mà điểm judge không phân biệt
được. Nếu full alignment rate và composite của harness cao hơn các đối sánh và khoảng
cách nới rộng khi mô hình nhỏ đi, điều đó ủng hộ trực tiếp cơ chế trung tâm: alignment
là thuộc tính suy giảm sớm trên mô hình nhỏ, và việc lưu nó thành trạng thái tường
minh giúp giữ lại thuộc tính này. Ngược lại, nếu các chỉ số alignment của harness
không tách khỏi các đối sánh dù điểm judge tách, thì lợi thế của harness đến từ các
khía cạnh chất lượng khác ngoài alignment---một kết quả vẫn hữu ích nhưng đòi hỏi hạ
claim về cơ chế. Sự đối chiếu giữa các chỉ số alignment và điểm judge cũng cho biết
rubric tổng quát nhạy đến đâu với vi phạm alignment.
\TBD{đối chiếu alignment--judge khi có số thực}.

\paragraph{Alignment graph như artifact giải thích được.}
Ngoài vai trò trạng thái cho máy, alignment graph còn là một sản phẩm cho con người.
Toàn bộ node, cạnh và trạng thái kiểm chứng được kết xuất thành file JSON kèm biểu
diễn Graphviz (\cref{sec:aggregation}), nên chuyên gia thiết kế giảng dạy có thể trả
lời các câu hỏi kiểm toán cụ thể---mục tiêu nào chưa có assessment tương ứng, liên kết
nào chưa qua kiểm chứng, sửa đổi nào đã kích hoạt kiểm tra lại---mà không cần đọc lại
toàn bộ văn bản chương trình. Điều này định vị hệ thống là công cụ hỗ trợ ra quyết
định (decision support) trong quy trình educator-in-the-loop: con người giữ quyền rà
soát, chỉnh sửa và override, còn harness bảo đảm rằng các quan hệ cần rà soát là tường
minh và truy vết được. Chúng tôi chủ ý tránh các claim kiểu ``đúng by construction'':
cơ chế chỉ làm cho alignment kiểm tra được và sửa được, không bảo đảm nó đúng.

\paragraph{Triển khai có trách nhiệm.}
Một chương trình sinh tự động có thể sai theo cách khó nhận ra: mục tiêu hoặc
assessment không phù hợp vẫn đọc trôi chảy. Vì vậy, đầu ra của hệ thống cần được
chuyên gia phê duyệt trước khi dùng với người học, đặc biệt trong các ngữ cảnh có hệ
quả cao; artifact đồ thị giảm chi phí của việc rà soát này nhưng không thay thế nó.
Nội dung sinh cũng có thể mang thiên lệch văn hoá, ngôn ngữ hoặc rào cản tiếp cận thừa
hưởng từ mô hình nền; các cơ chế của harness không kiểm tra các thuộc tính này và
chúng cần quy trình rà soát riêng. Cuối cùng, động cơ về riêng tư và tự lưu trữ chỉ
được hiện thực hoá trọn vẹn khi mô hình và judge chạy on-premise; cấu hình thực nghiệm
hiện tại vẫn dùng API bên ngoài (\cref{sec:threats-limitations}).
